\documentclass[12pt,a4paper]{article}

\usepackage[T1]{fontenc}
\usepackage[utf8]{inputenc}
\usepackage[brazil]{babel}
\usepackage{lmodern}
\usepackage{geometry}
\usepackage{hyperref}
\usepackage{enumitem}
\usepackage{array}

\geometry{margin=2.5cm}
\hypersetup{
  colorlinks=true,
  linkcolor=blue,
  urlcolor=blue
}

\title{Plano de Aula: Dependência Funcional e Normalização}
\date{}

\begin{document}
\maketitle

\noindent
\textbf{Duração:} 2 horas (120 minutos)\\
\textbf{Público-alvo:} Iniciantes\\
\textbf{Recursos Necessários:} Quadro branco/negro, marcadores/giz, folhas/cartões e (opcional) laboratório com SGBD para validação.

\section{Objetivos da Aula}
\begin{itemize}[leftmargin=*,itemsep=0.25em]
  \item Compreender dependência funcional e sua relação com chaves.
  \item Identificar redundância e anomalias de inserção/atualização/exclusão.
  \item Aplicar normalização até 3FN (quando apropriado).
  \item Melhorar esquemas relacionais para facilitar manutenção e consultas.
\end{itemize}

\section{Metodologia}
Aprendizagem baseada em problemas: apresentar uma tabela ``bagunçada'' (com repetição) e provocar anomalias. Em grupos, os alunos identificam dependências funcionais e reestruturam em tabelas menores, justificando a decomposição.

\section{Cronograma e Conteúdo Detalhado (120 Minutos)}
\subsection*{Parte 1: Problema Motivador (15 minutos)}
\begin{itemize}[leftmargin=*,itemsep=0.35em]
  \item Mostrar uma tabela única que mistura dados (ex.: Pedido + Cliente + Produto).
  \item Pergunta: ``O que acontece se o telefone do cliente mudar?'' (anomalia de atualização).
  \item Analogia: copiar a mesma informação em vários papéis gera inconsistência.
\end{itemize}

\subsection*{Parte 2: Dependência Funcional (25 minutos)}
\begin{itemize}[leftmargin=*,itemsep=0.35em]
  \item Definição: X determina Y (X \textrightarrow\ Y).
  \item Chaves candidatas e chave primária.
  \item Exercícios rápidos: identificar FDs em exemplos do quadro.
\end{itemize}

\subsection*{Parte 3: Formas Normais (30 minutos)}
\begin{itemize}[leftmargin=*,itemsep=0.35em]
  \item \textbf{1FN:} atributos atômicos, sem grupos repetidos.
  \item \textbf{2FN:} eliminar dependência parcial (quando há chave composta).
  \item \textbf{3FN:} eliminar dependência transitiva.
\end{itemize}

\subsection*{Parte 4: Atividade em Grupo --- Normalizando um Caso (40 minutos)}
\begin{itemize}[leftmargin=*,itemsep=0.35em]
  \item \textbf{Entrada:} tabela única ``Matrícula'' (Aluno, Curso, Disciplina, Professor, Sala, Nota).
  \item \textbf{Tarefas:}
  \begin{itemize}[leftmargin=*,itemsep=0.25em]
    \item Listar dependências funcionais observadas.
    \item Identificar anomalias (inserção, atualização, exclusão).
    \item Decompor em tabelas normalizadas e indicar PK/FK.
  \end{itemize}
  \item \textbf{Entrega:} conjunto de tabelas resultantes + justificativa breve.
\end{itemize}

\subsection*{Parte 5: Encerramento e Próximos Passos (10 minutos)}
\begin{itemize}[leftmargin=*,itemsep=0.35em]
  \item Revisão: por que normalizar e quando fazer.
  \item Ponte: ``Próximas aulas: recursos do SGBD --- procedures, asserções, gatilhos e transações.''
\end{itemize}

\end{document}

